\documentclass[11pt,a4paper]{article}
\usepackage[utf8]{inputenc}
\usepackage[spanish]{babel}
\usepackage{graphicx}
\usepackage{geometry}
\geometry{top=2.5cm, bottom=2.5cm, left=2.5cm, right=2.5cm}
\usepackage{xcolor}
\usepackage{listings}
\usepackage{tikz}
\usetikzlibrary{shapes,arrows,positioning,calc}

% Configuración de colores para código Java
\definecolor{javared}{rgb}{0.6,0,0}
\definecolor{javagreen}{rgb}{0.25,0.5,0.35}
\definecolor{javapurple}{rgb}{0.5,0,0.35}
\definecolor{javadocblue}{rgb}{0.25,0.35,0.75}

\lstset{
  language=Java,
  basicstyle=\ttfamily\footnotesize,
  keywordstyle=\color{javapurple}\bfseries,
  stringstyle=\color{javared},
  commentstyle=\color{javagreen},
  morecomment=[s][\color{javadocblue}]{/**}{*/},
  numbers=left,
  numberstyle=\tiny\color{black},
  stepnumber=1,
  numbersep=10pt,
  tabsize=4,
  showspaces=false,
  showstringspaces=false,
  frame=single,
  rulecolor=\color{black!30},
  breaklines=true,
  breakatwhitespace=true,
  captionpos=b,
  backgroundcolor=\color{gray!5}
}

% Comandos TikZ para diagramas UML
\tikzset{
    class/.style={
        rectangle,
        draw=black,
        thick,
        fill=blue!5,
        text width=6cm,
        align=center,
        minimum height=1cm
    },
    interface/.style={
        rectangle,
        draw=black,
        thick,
        fill=green!5,
        text width=6cm,
        align=center,
        minimum height=1cm,
        rounded corners=2mm
    },
    line/.style={
        draw,
        thick,
        -latex
    },
    implement/.style={
        draw,
        thick,
        dashed,
        -latex
    }
}

\title{\textbf{Documentación Arquitectónica: Patrones GoF}\\ \Large Sistema ReportaTuCalle}
\author{Ingeniería de Software}
\date{\today}

\begin{document}

\maketitle

\begin{abstract}
Este documento técnico detalla los Patrones de Diseño Gang of Four (GoF) implementados estrictamente a mano en el sistema \textit{ReportaTuCalle}. No se incluyen patrones que sean subproductos de frameworks (como la Inyección de Dependencias de Spring), sino código nativo en Java que fomenta la Alta Cohesión, el Bajo Acoplamiento y respeta la Arquitectura Hexagonal.
\end{abstract}

\tableofcontents
\newpage

\section{Patrones Estructurales y de Comportamiento}

\subsection{1. Strategy Pattern (Estrategia)}
\textbf{Problema:} El cálculo de optimización de rutas para el Supervisor debía soportar TSP, Flujo Máximo y Árbol de Expansión Mínima. Un \texttt{switch} gigante viola el principio Abierto/Cerrado (OCP).

\textbf{Solución:} Se programó una interfaz \texttt{OptimizationStrategy} y tres implementaciones concretas.

\subsubsection*{Diagrama UML (Strategy)}
\begin{center}
\begin{tikzpicture}[node distance=2cm and 1cm]
    % Nodos
    \node[class] (context) {\textbf{RouteOptimizationService} \\ \hline - strategies: List$<$OptimizationStrategy$>$ \\ \hline + optimizeRoute(...) };
    \node[interface, below=of context] (strategy) {$<<$Interface$>>$ \\ \textbf{OptimizationStrategy} \\ \hline + supports(type): boolean \\ + optimize(start, dest): OptimizedRoute };
    
    \node[class, text width=4cm, below left=of strategy, xshift=-1cm] (tsp) {\textbf{RoutingStrategyImpl} \\ \hline + optimize() };
    \node[class, text width=4cm, below=of strategy] (flow) {\textbf{FlowStrategyImpl} \\ \hline + optimize() };
    \node[class, text width=4cm, below right=of strategy, xshift=1cm] (mst) {\textbf{ConnectivityStrategyImpl} \\ \hline + optimize() };

    % Relaciones
    \path[line, draw, thick, -{Diamond[open]}] (context.south) -- (strategy.north);
    \path[implement] (tsp.north) -- ++(0,1) -| (strategy.south);
    \path[implement] (flow.north) -- (strategy.south);
    \path[implement] (mst.north) -- ++(0,1) -| (strategy.south);
\end{tikzpicture}
\end{center}

\subsubsection*{Fragmento de Código Extraído}
\begin{lstlisting}[caption={Interfaz OptimizationStrategy pura}]
public interface OptimizationStrategy {
    boolean supports(AlgorithmType type);
    Optional<OptimizedRoute> optimize(Coordinate startPoint, List<Coordinate> destinations);
}
\end{lstlisting}


\subsection{2. Chain of Responsibility (Cadena de Responsabilidad)}
\textbf{Problema:} La validación de un Reporte (Spam, Insultos, Ubicación) en el \texttt{ReportService} ensucia la lógica de negocio principal con múltiples validaciones anidadas.

\textbf{Solución:} Una cadena de validadores. Cada eslabón decide si procesa el reporte o lanza un error.

\subsubsection*{Fragmento de Código Extraído}
\begin{lstlisting}[caption={Implementación Base del Validador}]
public abstract class ReportValidationHandler {
    private ReportValidationHandler next;

    public void setNext(ReportValidationHandler next) {
        this.next = next;
    }

    public void validate(CreateReportRequest request) {
        doValidate(request); // Metodo Template
        if (next != null) {
            next.validate(request);
        }
    }
    protected abstract void doValidate(CreateReportRequest request);
}
\end{lstlisting}

\newpage

\section{Patrones Creacionales}

\subsection{3. Singleton (Puro en Java)}
\textbf{Problema:} Necesitamos guardar configuraciones estáticas en memoria (como límites de distancia o configuración global) de manera segura en un entorno multihilo, sin depender de la anotación \texttt{@Service} de Spring.

\textbf{Solución:} Implementación manual del Patrón Singleton utilizando el enfoque \textit{Double-Checked Locking} y \texttt{volatile}.

\subsubsection*{Diagrama UML (Singleton)}
\begin{center}
\begin{tikzpicture}
    \node[class] (singleton) {\textbf{AppConfigurationManager} \\ \hline - static instance: AppConfigurationManager \\ - configMap: Map$<$String, String$>$ \\ \hline - AppConfigurationManager() \\ + static getInstance(): AppConfigurationManager \\ + getConfig(key): String };
    
    \draw[line] (singleton.east) -- ++(1,0) |- ++(0,1.5) -| (singleton.north);
\end{tikzpicture}
\end{center}

\subsubsection*{Fragmento de Código Extraído}
\begin{lstlisting}[caption={AppConfigurationManager.java (Thread-Safe)}]
public final class AppConfigurationManager {
    private static volatile AppConfigurationManager instance;
    private final Map<String, String> configMap;

    private AppConfigurationManager() {
        this.configMap = new ConcurrentHashMap<>();
    }

    public static AppConfigurationManager getInstance() {
        if (instance == null) {
            synchronized (AppConfigurationManager.class) {
                if (instance == null) {
                    instance = new AppConfigurationManager();
                }
            }
        }
        return instance;
    }
}
\end{lstlisting}


\section{Conclusión de la Arquitectura}
Se aplicó estrictamente el principio \textbf{YAGNI} al programar para interfaces (sólo se usaron interfaces donde el comportamiento variará, como en las \texttt{OptimizationStrategy} y los Adaptadores de la Arquitectura Hexagonal). Se ha garantizado \textbf{Alta Cohesión} modularizando por entidad y \textbf{Bajo Acoplamiento} inyectando comportamiento vía Decoradores y Observers.

\end{document}
